\section{Введение}
\label{sec:introduction}
Программирующий агент узнаёт задачу через действия. Команда обнаруживает зависимость, тест опровергает исправление, сообщение компилятора связывает ошибку с конкретной правкой. Последующие решения зависят от того, какие сведения об этих действиях остались доступны. Если история взаимодействия растёт на каждом шаге, вместе с ней растёт и объём входа модели. Если её сжать, приходится решать, чем заменить удалённые наблюдения.

SKILL.state предлагает хранить изменяемое структурированное состояние вместо полного транскрипта~\citep{skillstate2026}. Привлекательность этого решения состоит в отделении текущей рабочей памяти от накопленной истории. При ограниченном размере состояния контекст не обязан расти с каждым действием. Но ограничение размера представления само по себе не ограничивает объём информации, необходимой для решения задачи. Краткая запись о неудаче может не сохранить сообщение об ошибке; список изменённых файлов не сообщает, что именно было изменено и почему.

Из этого возникает инженерный вопрос: какую работу по ведению памяти следует поручить модели? Гипотезы и планы требуют интерпретации. Выполненная команда и её наблюдённый результат уже известны среде исполнения. Повторная генерация этих сведений моделью расходует выходные токены и делает их сохранность зависимой от обновлений семантического состояния. Разделение обязанностей позволяет хранить наблюдения автоматически, оставляя модели рассуждение об их значении. Мы называем такое представление гибридным состоянием.

В работе проверяются возможности этого разделения на задачах программирующего агента. Сначала оценивается цена замены истории состоянием модели. Затем сравнивается состав автоматической памяти: от ссылок на файлы до записей изменений, команд, диагностик и результатов тестов. Следующие сравнения проверяют, устраняют ли оставшееся отставание недавняя история, адаптация схемы или другой протокол обновления. Качество рассматривается вместе с ресурсами, поскольку меньший промпт может сочетаться с большей генерацией и более высокой стоимостью успешного выполнения.

Результаты дают различающиеся ответы на эти вопросы. Компактное состояние действительно уменьшает пиковый контекст, но без истории теряет успешность. В поисковом сравнении более содержательные записи среды исполнения заметно лучше бедного списка файлов; дальнейшее обогащение не даёт последовательного улучшения. В повторных измерениях наиболее выраженное восстановление успешности связано с недавней историей. Испытанные схема и inline-протокол не обеспечивают аналогичного результата. При этом экономический эффект сохранения истории различается между реализациями.

Вклад работы состоит в эмпирическом разделении решений, которые легко объединить под названием <<память агента>>: кто записывает сведения, что входит в состояние, какие взаимодействия остаются в исходном виде и как оплачивается дальнейшая работа. Это исследование переноса идеи SKILL.state в другой агентный стек, а не повторение исходных экспериментов. Его результаты характеризуют испытанные реализации и уточняют, какие изменения частично компенсируют их ограничения.
